\subsubsection{Parametrization of the Flux-Surfaces}

In this case, we have $\psi=\psi\left(  r\right)  $ and therefore it is easier
to work in the $\left(  r,\theta,\phi\right)  $ coordinate to account for the
symmetry in the problem. The normalized radius is
\begin{equation}
\rho\left(  \psi\right)  =\frac{r}{a}%
\end{equation}

We have now
\begin{equation}
\widehat{\psi}=\widehat{r}%
\end{equation}
so that (\ref{gijrtheta})
\begin{equation}
\sqrt{g_{22}}=r\label{g22circ}%
\end{equation}

\subsubsection{Magnetic Field}

The toroidal field is
\begin{equation}
B_{T}\left(  \rho,\theta\right)  =\frac{\left|  I\left(  \rho\right)  \right|
}{R\left(  \rho,\theta\right)  }%
\end{equation}
and the poloidal field is
\begin{equation}
B_{P}\left(  \rho,\theta\right)  =\frac{\left|  \nabla\psi\left(  \rho\right)
\right|  }{R\left(  \rho,\theta\right)  }\label{Bpcirc}%
\end{equation}
where
\begin{equation}
\left|  \nabla\psi\left(  \rho\right)  \right|  =\frac{1}{a}\left|
\frac{d\psi\left(  \rho\right)  }{d\rho}\right|
\end{equation}
is now only a function of $\rho$.

The total field is then
\begin{equation}
B\left(  \rho,\theta\right)  =\frac{\sqrt{I^{2}\left(  \rho\right)  +\left|
\nabla\psi\left(  \rho\right)  \right|  ^{2}}}{R\left(  \rho,\theta\right)  }%
\end{equation}
and can be written as
\begin{equation}
B\left(  \rho,\theta\right)  =B_{0}\left(  \rho\right)  \frac{R_{0}}{R\left(
\rho,\theta\right)  }\label{Bcirc}%
\end{equation}
with
\begin{equation}
B_{0}\left(  \rho\right)  =\frac{\sqrt{I^{2}\left(  \rho\right)  +\left|
\nabla\psi\left(  \rho\right)  \right|  ^{2}}}{R_{0}}%
\end{equation}

Consequently, we ratio of magnetic fields is (\ref{Psidef})%
\begin{equation}
\Psi\left(  \rho,\theta\right)  =\frac{R_{0}}{R\left(  \rho,\theta\right)
}\label{Psidefcirc}%
\end{equation}

We have the ratio
\begin{align}
\frac{B}{B_{P}}  & =\frac{\sqrt{I^{2}\left(  \rho\right)  +\left|  \nabla
\psi\left(  \rho\right)  \right|  ^{2}}}{\left|  \nabla\psi\left(
\rho\right)  \right|  }\label{BsBpcirc}\\
& =\sqrt{1+\frac{I^{2}\left(  \rho\right)  }{\left|  \nabla\psi\left(
\rho\right)  \right|  ^{2}}}%
\end{align}

\subsubsection{Safety factor}

The safety factor becomes (\ref{qfactor})
\begin{equation}
q\left(  \psi\right)  =\int_{0}^{2\pi}\frac{d\theta}{2\pi}\frac{r}{R}%
\frac{B_{T}}{B_{P}}=\frac{r}{R_{p}}\frac{B_{T}}{B_{P}}\int_{0}^{2\pi}%
\frac{d\theta}{2\pi}\frac{R_{p}}{R}%
\end{equation}

The averaged value of $R_{p}/R$ is evaluated in (\ref{sstar}). It gives
\begin{equation}
\int_{0}^{2\pi}\frac{d\theta}{2\pi}\frac{R_{p}}{R}=\frac{1}{\sqrt{1-\left(
r/R_{p}\right)  ^{2}}}%
\end{equation}
so that
\begin{equation}
q\left(  \psi\right)  =\frac{1}{\sqrt{1-\left(  r/R_{p}\right)  ^{2}}}\frac
{r}{R_{p}}\frac{B_{T}}{B_{P}}\label{qcirc}%
\end{equation}

Note that in the factor $\sqrt{1-\left(  r/R_{p}\right)  ^{2}}$ is usually
neglected, which is valid only in the large aspect ratio approximation.

\subsubsection{Particle Motion}

Defining
\begin{equation}
\epsilon=\frac{r}{R_{p}}%
\end{equation}
and using (\ref{Rdef}),%
\begin{equation}
R\left(  \rho,\theta\right)  =R_{p}+r\cos\theta
\end{equation}

The position of minimum $B$-field is
\begin{equation}
\theta_{0}=0
\end{equation}
and we find
\begin{align}
R_{\min}\left(  r\right)    & =R_{p}-r=R_{p}\left(  1-\epsilon\right)  \\
R_{\max}\left(  r\right)    & =R_{p}+r=R_{p}\left(  1+\epsilon\right)
=R_{0}\left(  r\right)  \nonumber
\end{align}
and (\ref{Psidefcirc})%
\begin{equation}
\Psi\left(  \rho,\theta\right)  =\frac{1+\epsilon}{1+\epsilon\cos\theta
}\label{Psidefcirc2}%
\end{equation}
so that, using (\ref{Bcirc})
\begin{equation}
B_{\max}\left(  r\right)  =B_{0}\left(  r\right)  \frac{1+\epsilon}%
{1-\epsilon}\nonumber
\end{equation}
and, using (\ref{xi0t2})
\begin{equation}
\xi_{0T}^{2}\left(  r\right)  =\frac{2\epsilon}{1+\epsilon}%
\end{equation}

The pitch-angle cosine is then given by (\ref{xi}) (\ref{Psidefcirc2})
\begin{equation}
\xi\left(  r,\theta,\xi_{0}\right)  =\text{sign}\left(  \xi_{0}\right)
\sqrt{1-\frac{1+\epsilon}{1+\epsilon\cos\theta}\left(  1-\xi_{0}^{2}\right)
}\label{xicirc}%
\end{equation}

The turning angles are determined by
\begin{equation}
B\left(  r,\theta_{T}\right)  =B_{b}\left(  r,\xi_{0}\right)  =\frac
{B_{0}\left(  r\right)  }{1-\xi_{0}^{2}}%
\end{equation}
which gives, using (\ref{Psidefcirc2}),
\begin{equation}
B_{0}\left(  r\right)  \frac{1+\epsilon}{1+\epsilon\cos\theta_{T}}=\frac
{B_{0}\left(  r\right)  }{1-\xi_{0}^{2}}%
\end{equation}
and then
\begin{align}
\xi_{0}^{2}  & =1-\frac{1+\epsilon\cos\theta_{T}}{1+\epsilon}\nonumber\\
& =\frac{\epsilon\left(  1-\cos\theta_{T}\right)  }{1+\epsilon}%
\end{align}
so that
\begin{equation}
\theta_{T}=\arccos\left[  1-\frac{2\xi_{0}^{2}}{\xi_{0T}^{2}}\right]
\end{equation}
and by symmetry
\begin{align}
\theta_{T\min}  & =-\theta_{T}\\
\theta_{T\max}  & =\theta_{T}\nonumber
\end{align}

\subsubsection{Bounce Time}

Using (\ref{g22circ}), the normalized bounce time reduces to
\begin{equation}
\lambda\left(  \rho,\xi_{0}\right)  =\frac{1}{\widetilde{q}\left(
\psi\right)  }\frac{r}{R_{p}}\int_{\theta_{\min}}^{\theta_{\max}}\frac
{d\theta}{2\pi}\frac{\xi_{0}}{\xi}\frac{B}{B_{P}}%
\end{equation}
with, using (\ref{qtilde})
\begin{equation}
\widetilde{q}\left(  \psi\right)  =\frac{r}{R_{p}}\int_{0}^{2\pi}\frac
{d\theta}{2\pi}\frac{B}{B_{P}}\label{qtilde2}%
\end{equation}

Because $B/B_{P}$ only a function of $\rho$, as seen in (\ref{BsBpcirc}), and
can be taken out of the integrals, we get finally%

\begin{equation}
\lambda\left(  \rho,\xi_{0}\right)  =\int_{\theta_{\min}}^{\theta_{\max}}%
\frac{d\theta}{2\pi}\frac{\xi_{0}}{\xi}\label{lambdacirc}%
\end{equation}

This integral can be performed analytically in a series expansion and are
calculated in (\ref{lambdacalcl}). Note that in the case where $B_{T}\gg
B_{P}$ and in the large aspect ratio approximation, we have $\widetilde
{q}\left(  \psi\right)  \rightarrow q\left(  \psi\right)  $, which explains
the notations.
